\documentclass[instructions.tex]{subfiles}
\begin{document}
 
\chapter{De l'Extrême-Onction.}
 
Ce sacrement, dit le concile de Trente, est comme la consommation de la pénitence et de la vie chrétienne. Le Sauveur, ayant pourvu au salut des fidèles pendant leur vie, n'a pas voulu qu'ils fussent sans secours à la mort. L'apôtre saint Jacques, le parent du Seigneur, nous a fait connaître l'Institution de ce sacrement et ses effets.

\primo Ce sacrement augmente la grace sanctifiante, achève de purifier l'âme, remet les péchés qui ne sont pas encore expiés, supplée aux pénitences que nous n'avons pas faites, plus ou moins, selon les dispositions de ceux qui le reçoivent; il fortifie contre les assaults des démons, qui attaquent les mourans avec plus de fureur, voyant qu'il leur reste peu de temps\footnote{Diabolus...... habens iram magnam, sciens quia modicum tempus habet, Ap. 12. 12.}; Ce sacrement console les mourans, adoucit leurs peines, les rassure contre les inquiétudes de la mort, et contre l'excessive frayeur des jugemens de Dieu; ranime leur confiance et leur amour envers Dieu; il remet enfin le sceau à la prédestination par une mort sainte. Combien seraient morts en reprouves sans ce remède salutaire! combien de fidèles craignant Dieu, qui d'abord après leur mort vont au ciel ou demeurent peu de temps en purgatoire, qui resteraient long-temps dans les tourmens, s'ils n'avaient pas reçu ce sacrement des mourans!

Sacrement si efficace, qu'il peut, selon la doctrine du saint concile, rétablir la santé quand il est expédient pour le salut, pourvu, dit saint Thomas, que le malade n'y mette pas obstacle. Plusieurs ont en effet recouvré leur santé, après avoir reçu l'huile sainte. On verrait bien plus souvent les effets de ce sacrement, et l'efficacité des promesses de Jésus-Christ, si on le recevait avec une vive foi. C'est donc avec raison que l'Eglise a frappé d'anathême ceux qui disent qu'un chrétien ne pèche point en le méprisant: Le mépris d'un si grand sacrement, dit le saint concile, est un crime énorme, une outrage contre le Saint-Esprit\footnote{Sess. 24. c. 3.}.

\secundo Quant à la préparation pour le recevoir, il faut deux choses, le temps et les dispositions. 1. Le temps. Ce sacrement ne doit être reçu que dans la maladie, lorsqu'on est dans un prochain danger de la mort. On ne doit pas cependant attendre qu'on soit à l'extrémité: ce serait tenter Dieu d'espérer la santé du corps par la vertu de ce sacrement, lorsqu'on attend pour le recevoir, que tout soit désespéré. Dieu n'est pas obligé de nous tirer par miracle, d'une extrémité où notre négligence nous a réduits, l'âme elle-même ne recevrait pas ses effets si abondans, si l'on attendait de recevoir ce sacrement quand on n'a plus ni connaissance ni sentiment. 2. Les dispositions. L'extrême-onction étant la consommation de la pénitence, il faut être pénitent, confessé, s'il se peut, pour le recevoir. Le secours du sacre viatique, le désir de s'unir à Dieu et de l'aimer, un vif repentir de ses fautes, une pleine confiance aux mérites du Sauveur: voilà ce qu'on espere la miséricorde par la vertu de l'onction sainte, sont les dispositions avec lesquelles ce sacrement produit ses effets dans leur plénitude.

C'est donc un piège du démon et une grande faiblesse de craindre ce sacrement; loin d'avancer la mort et de nuire au malade, il peut guérir le corps et sanctifier l'âme. On se met au hasard d'en être privé lorsqu'on diffère trop; et l'on s'expose au danger de n'en pas ressentir les effets, quand on le reçoit avec répugnance ou par pure bienséance.

Ceux qui ont soin des malades doivent les avertir de bonne heure, de faire leur testament. Faut-il, pour régler les affaires de la famille, attendre à un temps où le malade n'a plus quelques momens pour assurer son salut et régler sa conscience, qui est peut-être encore dans le desordre? Qu'on laisse périr tout le reste, s'il le faut; mais qu'on ne laisse pas périr son ame.

On manque au devoir essentiel de la charité lorsqu'on, par respect humain ou par intérêt, on lui cache le danger où il est. On lui fait tort, si, dans la crainte de l'alarmer, et par des vues humaines, on l'expose à mourir sans savoir qu'il meurt. Pourquoi lui ôter le mérite et la consolation d'offrir à Dieu, le sacrifice de sa vie? Combien de péchés n'effacera-t-il pas ce sacrifice volontaire? Pourquoi l'exposer à recevoir sans connaissance, et peut-être sans fruit, un sacrement si consolant? Quel préjudice ne porte-t-on pas à un mourant par ce silence pernicieux!

On doit exhorter un mourant avec des paroles tendres et affectueuses, en peu de mots, les répéter de temps en temps, pas trop souvent, ni d'une voix trop élevée, qui fatiguerait le malade. Qu'on n'oublie pas l'eau bénite dont l'aspersion peut mettre les démons en fuite\footnote{Ad abigendos dæmones, Orat. Eccl.}.

\section{Résolutions}

\primo Puisque ce sacrement produit tant et de si heureux effets, je demanderai qu'on me l'administre, dès que je serai atteint griévement d'une maladie dangereuse. 

\secundo Je m'y préparerai avec soin par la confession, accompagnée de grands sentimens de confiance et foi. 

\tertio J'aurai soin aussi, autant que je le pourrai, que mes proches, mes amis, en général mes connaissances et tous ceux qui dépendront de moi, ne décèdent pas sans ce puissant secours. 
\prayer{Soyez béni, ô mon Dieu! de ce que vous avez bien voulu nous donner un moyen si consolant et si efficace pour nous communiquer vos grâces, dans une circonstance où nous en avons tant besoin.}
 
\end{document}
